\chapter{Introduction}
\label{ch:intro}

\section{Background}

Almost every large language model (LLM) in use today is built on the
Transformer architecture introduced by Vaswani et al.\ in
2017~\cite{vaswani2017attention}. Its central operation, self-attention, lets
every token in the input look directly at every other token, which gives the
model an exceptionally flexible way of routing information across a sequence.
That flexibility has a price. The amount of computation in an attention layer
grows quadratically with the length of the input, and during text generation
the model must keep the keys and values of all previous tokens in a
\emph{key--value (KV) cache} whose size grows with every token that is added.
For a chat message of a few hundred tokens this cost is negligible. For the
workloads that are now becoming common---an AI agent that reads an entire code
repository, works through a long tool-calling session, or analyses a
several-hundred-page document---it becomes the dominant cost of running the
model.

Two independent lines of research have attacked this cost from different
directions. The first is the \emph{state-space model} (SSM). Mamba~\cite{gu2023mamba}
and its successor Mamba-2~\cite{dao2024mamba2} read a sequence one token at a
time and summarise everything seen so far in a state of fixed size, so that
time grows only linearly with the sequence length and the memory needed during
generation does not grow at all. The second is the \emph{Mixture-of-Experts}
(MoE) layer~\cite{shazeer2017moe,lepikhin2020gshard}, which splits the
feed-forward part of a model into many small ``experts'' and activates only a
handful of them for each token, so that a model can hold a very large number of
parameters while spending the computation of a much smaller one.

\NS, released by NVIDIA in April 2026~\cite{nvidia2026super}, combines both
ideas in a single open model. It has 120.6~billion parameters in total, of
which only 12.7~billion are active in any forward pass. Most of its 88 layers
are Mamba-2 and MoE layers; only eight are self-attention layers. The model
accepts contexts of up to one million tokens and, according to its technical
report, achieves accuracy comparable to GPT-OSS-120B~\cite{openai2025gptoss}
and Qwen3.5-122B while delivering up to 2.2$\times$ and 7.5$\times$ higher
inference throughput than those models respectively.

\section{Motivation}

The motivation for studying this model is that it is the clearest example so
far of a production-scale LLM that is \emph{not} a pure Transformer. Earlier
hybrid models were either small or research prototypes. \NS\ shows that the
hybrid approach scales to more than one hundred billion parameters and to
million-token contexts while remaining competitive with the strongest open
Transformer models of the same size. Understanding how it works therefore
answers a broader question: is the attention-only design that has dominated
since 2017 still the best foundation for large models, or is the field moving
towards mixed architectures in which attention is used sparingly?

A second motivation is that the model is fully open. Its weights, a large part
of its training data and its training recipe are published, and the technical
report describes the engineering in unusual detail---including problems that
appeared along the way, such as the numerical drift caused by quantizing the
Mamba state. This makes it possible to study not only the architecture but
also what it takes to train and deploy it.

\section{Objectives of the Seminar}

The objectives of this seminar are:
\begin{itemize}
  \item to understand why Transformers become slow and memory-heavy on long
        inputs;
  \item to explore state-space models, in particular Mamba and Mamba-2, which
        were proposed to solve that problem;
  \item to study how \NS\ interleaves Mamba-2 and attention layers to form a
        hybrid model;
  \item to examine how Mixture-of-Experts, and specifically LatentMoE, lets the
        model stay large while running cheaply;
  \item to review how the model is trained, quantized and evaluated; and
  \item to highlight what this design suggests about the future of efficient
        large-scale AI.
\end{itemize}

\section{Scope}

The report is a study of published work. Its primary source is the \NS\
technical report~\cite{nvidia2026super}; the supporting papers on Mamba,
Mamba-2, LatentMoE, multi-token prediction and NVFP4 training are used to
explain the building blocks in more depth. All benchmark numbers quoted in this
report are those published by the respective authors; no new experiments were
run. Where the report contains estimates worked out by the author from
published configuration values, this is stated explicitly.

\section{Organisation of the Report}

Chapter~\ref{ch:survey} surveys the literature that leads up to \NS.
Chapter~\ref{ch:background} gives the technical background on Transformers and
the KV cache, state-space models, and Mixture-of-Experts.
Chapter~\ref{ch:arch} describes the architecture of \NS: the hybrid layout,
LatentMoE and multi-token prediction. Chapter~\ref{ch:training} covers
pretraining in the 4-bit NVFP4 format, the training data, long-context
extension, supervised fine-tuning and reinforcement learning.
Chapter~\ref{ch:quant} explains how the model is quantized for deployment and
how the Mamba state is handled. Chapter~\ref{ch:results} presents the results,
Chapter~\ref{ch:apps} the applications, and
Chapter~\ref{ch:adv} the advantages and limitations.
Chapter~\ref{ch:conclusion} concludes the report and outlines future scope.
